% Generated by writing/make_appendix_assets.py; do not edit.
% source: processing/Processing/output/paper/tables/ideological_universe_comparison.csv
\begin{table}[!htbp]
\centering\small
\caption{Minimal winning coalitions by ideological domain and party universe}
\label{tab:one-gap-summary}
\label{tab:ideological-universe-comparison}
\setlength{\tabcolsep}{6pt}
\renewcommand{\arraystretch}{1.12}
\begin{tabular}{@{}lcrrrr@{}}
\toprule
 & & \multicolumn{2}{c}{Seat-winning parties} & \multicolumn{2}{c}{All parties} \\
\cmidrule(lr){3-4}\cmidrule(l){5-6}
Election & \(k\) & Minimal majorities & Inversions & Minimal majorities & Inversions \\
\midrule
2014 & 0 & 8 & 4 & 8 & 4\\
2014 & 1 & 87 & 46 & 88 & 43\\
\addlinespace[3pt]
2018 & 0 & 8 & 1 & 8 & 0\\
2018 & 1 & 108 & 42 & 118 & 24\\
\addlinespace[3pt]
2022 & 0 & 4 & 2 & 4 & 2\\
2022 & 1 & 39 & 12 & 53 & 17\\
\bottomrule
\end{tabular}
\par\vspace{4pt}\begin{minipage}{\textwidth}\footnotesize\textit{Note:} Minimality is recomputed for each \(k\) and universe. The primary \(k=1\) rows sum to 234 minimal majorities and 100 inversions. The \(k=0\) and \(k=1\) rows are overlapping domains and are not pooled. Every vote share uses all valid federal-deputy votes in \(V\).
\end{minipage}
\end{table}
